% Amateur Radio Operating Reference — common Q codes first edition
% Compile with XeLaTeX or LuaLaTeX for Arial, or with pdfLaTeX for Helvetica.
\documentclass[11pt,letterpaper]{article}
\usepackage[margin=0.68in,top=0.64in,bottom=0.62in,headheight=14pt,headsep=12pt,footskip=24pt]{geometry}
\usepackage{iftex}
\ifPDFTeX
  \usepackage[T1]{fontenc}
  \usepackage[utf8]{inputenc}
  \usepackage[scaled=0.95]{helvet}
  \renewcommand{\familydefault}{\sfdefault}
\else
  \usepackage{fontspec}
  \IfFontExistsTF{Arial}{\setmainfont{Arial}}{\setmainfont{TeX Gyre Heros}}
\fi
\usepackage[table]{xcolor}
\usepackage{array,tabularx,fancyhdr}
\usepackage{ragged2e}
\usepackage[hidelinks]{hyperref}
\definecolor{Ink}{HTML}{111827}
\definecolor{TitleInk}{HTML}{102A43}
\definecolor{Accent}{HTML}{00549A}
\definecolor{Secondary}{HTML}{384D61}
\definecolor{TableHead}{HTML}{DFEDF5}
\definecolor{RowTint}{HTML}{F1F5F8}
\definecolor{Grid}{HTML}{9BAAB4}
\color{Ink}
\pagestyle{fancy}
\fancyhf{}
\fancyhead[R]{\fontsize{8}{9}\selectfont\bfseries\color{Secondary} AMATEUR RADIO / OPERATING REFERENCE}
\fancyfoot[L]{\fontsize{8}{9}\selectfont\color{Secondary} Source: arrl.org/operating-aids \quad\textbullet\quad Page \thepage}
\renewcommand{\headrulewidth}{0pt}
\renewcommand{\footrulewidth}{0pt}
\setlength{\parindent}{0pt}
\setlength{\parskip}{3pt}
\setlength{\tabcolsep}{4pt}
\renewcommand{\arraystretch}{1.0}
\arrayrulecolor{Grid}
\newcolumntype{Y}[1]{>{\hsize=#1\hsize\linewidth=\hsize\RaggedRight\arraybackslash}X}
\newcommand{\pagetitle}[1]{{\color{TitleInk}\fontsize{21}{25}\selectfont\bfseries #1}\par\vspace{2pt}{\color{Accent}\hrule height 0.5pt}\vspace{4pt}}
\newcommand{\subtitle}[1]{{\color{Secondary}\itshape #1}\par\vspace{4pt}}
\newcommand{\sectiontitle}[1]{\par\vspace{5pt}{\color{Accent}\large\bfseries #1}\par\vspace{2pt}}
\begin{document}
\pagetitle{Phonetics, signals \& UTC}
\subtitle{Quick reference for clear identification, signal reports and time conversion}
\sectiontitle{Phonetic alphabet}
{\fontsize{10}{12.1}\selectfont
\rowcolors{2}{white}{RowTint}
\begin{tabularx}{\linewidth}{|Y{0.43}|Y{0.79}|Y{1.78}|Y{0.43}|Y{0.83}|Y{1.74}|}
\hline
\rowcolor{TableHead}
\textbf{Ltr.} & \textbf{Word} & \textbf{Pronunciation} & \textbf{Ltr.} & \textbf{Word} & \textbf{Pronunciation} \\ \hline
\textbf{A} & Alpha & AL fah & \textbf{N} & November & no VEM ber \\ \hline
\textbf{B} & Bravo & BRAH voh & \textbf{O} & Oscar & OSS cah \\ \hline
\textbf{C} & Charlie & CHAR lee & \textbf{P} & Papa & pah PAH \\ \hline
\textbf{D} & Delta & DELL tah & \textbf{Q} & Quebec & keh BECK \\ \hline
\textbf{E} & Echo & ECK oh & \textbf{R} & Romeo & ROW me oh \\ \hline
\textbf{F} & Foxtrot & FOKS trot & \textbf{S} & Sierra & see AIR rah \\ \hline
\textbf{G} & Golf & GOLF & \textbf{T} & Tango & TANG go \\ \hline
\textbf{H} & Hotel & hoh TELL & \textbf{U} & Uniform & YOU nee form \\ \hline
\textbf{I} & India & IN dee ah & \textbf{V} & Victor & VIK tah \\ \hline
\textbf{J} & Juliet & JEW lee ETT & \textbf{W} & Whiskey & WISS key \\ \hline
\textbf{K} & Kilo & KEY loh & \textbf{X} & X-ray & ECKS ray \\ \hline
\textbf{L} & Lima & LEE mah & \textbf{Y} & Yankee & YANG kee \\ \hline
\textbf{M} & Mike & MIKE & \textbf{Z} & Zulu & ZOO loo \\ \hline
\end{tabularx}
}\par
\sectiontitle{Signal reporting: RST}
{\fontsize{10}{12.1}\selectfont
\rowcolors{2}{white}{RowTint}
\begin{tabularx}{\linewidth}{|Y{0.78}|Y{0.35}|Y{1.87}|}
\hline
\rowcolor{TableHead}
\textbf{Component} & \textbf{Scale} & \textbf{Meaning} \\ \hline
R --- Readability & 1--5 & 5 is best. \\ \hline
S --- Signal strength & 1--9 & Strength of the received signal. \\ \hline
T --- Tone & 1--9 & Quality of a CW note; modern equipment is rarely below 9. \\ \hline
\end{tabularx}
}\par
Example: A signal 30 dB over S-9 can be reported as ``five nine plus thirty'' or ``thirty over nine.''\par
\sectiontitle{Convert local time to UTC}
Add the hours below to local time and use 24-hour notation. If the result passes midnight, advance the UTC date.\par
{\fontsize{10}{12.1}\selectfont
\rowcolors{2}{white}{RowTint}
\begin{tabularx}{\linewidth}{|Y{1}|Y{1}|Y{1}|}
\hline
\rowcolor{TableHead}
\textbf{Time zone} & \textbf{Standard time} & \textbf{Daylight saving time} \\ \hline
Atlantic & +4 hours & +3 hours \\ \hline
Eastern & +5 hours & +4 hours \\ \hline
Central & +6 hours & +5 hours \\ \hline
Mountain & +7 hours & +6 hours \\ \hline
Pacific & +8 hours & +7 hours \\ \hline
\end{tabularx}
}\par
Examples: 10:00 AM EST = 1500 UTC.  7:00 PM PST on Dec 31 = 0300 UTC on Jan 1.\par
\newpage
\pagetitle{Q signals \& procedure}
\subtitle{Read across: commonly used codes first, then others alphabetically. Usage varies by mode. A question mark asks a question; without it, a code gives a statement or instruction.}
{\fontsize{10}{12.1}\selectfont
\rowcolors{2}{white}{RowTint}
\begin{tabularx}{\linewidth}{|Y{0.35}|Y{1.65}|Y{0.35}|Y{1.65}|}
\hline
\rowcolor{TableHead}
\textbf{Code} & \textbf{Question / statement} & \textbf{Code} & \textbf{Question / statement} \\ \hline
\textbf{QSL} & Can you acknowledge receipt? & \textbf{QSO} & Can you communicate with \ldots{} direct? \\ \hline
\textbf{QTH} & What is your location? & \textbf{QRZ} & Who is calling me? \\ \hline
\textbf{QRM} & Is my transmission being interfered with? & \textbf{QRN} & Are you troubled by static? \\ \hline
\textbf{QSY} & Shall I change frequency? & \textbf{QSB} & Are my signals fading? \\ \hline
\textbf{QRP} & Shall I decrease transmitter power? & \textbf{QRT} & Shall I stop sending? \\ \hline
\textbf{QRL} & Are you busy? & \textbf{QRS} & Shall I send slower? \\ \hline
\textbf{QRV} & Are you ready? & \textbf{QRX} & When will you call again? \\ \hline
\textbf{QRO} & Shall I increase transmitter power? & \textbf{QRQ} & Shall I send faster? \\ \hline
\textbf{QSK} & Can you work break-in? & \textbf{QRU} & Have you anything for me? (Answer in negative.) \\ \hline
\textbf{QST} & General call: message to all amateurs / ARRL members. & \textbf{QRA} & What is the name of your station? \\ \hline
\textbf{QRB} & How far are you from my station? & \textbf{QRG} & What is my exact frequency? \\ \hline
\textbf{QRH} & Does my frequency vary? & \textbf{QRI} & How is the tone of my transmission? (1--3) \\ \hline
\textbf{QRJ} & Are you receiving me badly? & \textbf{QRK} & How intelligible are my signals? (1--5) \\ \hline
\textbf{QRW} & Shall I tell \ldots{} you are calling? & \textbf{QRY} & What is my turn? \\ \hline
\textbf{QSA} & What is my signal strength? (1--5) & \textbf{QSD} & Is my keying defective? \\ \hline
\textbf{QSG} & Shall I send \ldots{} messages at a time? & \textbf{QSM} & Shall I repeat the last message? \\ \hline
\textbf{QSN} & Did you hear me on \ldots{} kHz? & \textbf{QSP} & Will you relay to \ldots{}? \\ \hline
\textbf{QSU} & Shall I send or reply on this frequency? & \textbf{QSV} & Shall I send a series of Vs? \\ \hline
\textbf{QSW} & Will you send on this frequency? & \textbf{QSX} & Will you listen for \ldots{} on \ldots{}? \\ \hline
\textbf{QSZ} & Shall I send each word or group more than once? & \textbf{QTA} & Shall I cancel message number \ldots{}? \\ \hline
\textbf{QTB} & Do you agree with my word count? & \textbf{QTC} & How many messages have you to send? \\ \hline
\textbf{QTE} & What is my true bearing from you? & \textbf{QTR} & What is your time? \\ \hline
\textbf{QTU} & What are your station's operating hours? & \textbf{QUA} & Have you news of \ldots{}? \\ \hline
\end{tabularx}
}\par
Scales: QSA is 1--5 (RST signal strength is 1--9). QRI tone: 1 good, 2 variable, 3 bad.\par
\sectiontitle{Communications procedure}
{\fontsize{10}{12.1}\selectfont
\rowcolors{2}{white}{RowTint}
\begin{tabularx}{\linewidth}{|Y{1.4}|Y{0.4}|Y{1.8}|Y{0.4}|}
\hline
\rowcolor{TableHead}
\textbf{Voice} & \textbf{Code} & \textbf{Voice} & \textbf{Code} \\ \hline
Go ahead & \textbf{K} & Roger (received in full) & \textbf{R} \\ \hline
Over & \textbf{AR} & Clear & \textbf{SK} \\ \hline
Over (specific station) & \textbf{KN} & Closing the station & \textbf{CL} \\ \hline
Stand by or wait & \textbf{AS} &  &  \\ \hline
\end{tabularx}
}\par
\par
\newpage
\pagetitle{State \& province abbreviations}
\subtitle{Alphabetical reference \textbullet{} U.S. states and Canadian provinces and territories}
{\fontsize{10}{12.1}\selectfont
\rowcolors{2}{white}{RowTint}
\begin{tabularx}{\linewidth}{|Y{1.62}|Y{0.38}|Y{1.62}|Y{0.38}|}
\hline
\rowcolor{TableHead}
\textbf{State} & \textbf{Abbr.} & \textbf{State} & \textbf{Abbr.} \\ \hline
Alabama & \textbf{AL} & Montana & \textbf{MT} \\ \hline
Alaska & \textbf{AK} & Nebraska & \textbf{NE} \\ \hline
Arizona & \textbf{AZ} & Nevada & \textbf{NV} \\ \hline
Arkansas & \textbf{AR} & New Hampshire & \textbf{NH} \\ \hline
California & \textbf{CA} & New Jersey & \textbf{NJ} \\ \hline
Colorado & \textbf{CO} & New Mexico & \textbf{NM} \\ \hline
Connecticut & \textbf{CT} & New York & \textbf{NY} \\ \hline
Delaware & \textbf{DE} & North Carolina & \textbf{NC} \\ \hline
Florida & \textbf{FL} & North Dakota & \textbf{ND} \\ \hline
Georgia & \textbf{GA} & Ohio & \textbf{OH} \\ \hline
Hawaii & \textbf{HI} & Oklahoma & \textbf{OK} \\ \hline
Idaho & \textbf{ID} & Oregon & \textbf{OR} \\ \hline
Illinois & \textbf{IL} & Pennsylvania & \textbf{PA} \\ \hline
Indiana & \textbf{IN} & Rhode Island & \textbf{RI} \\ \hline
Iowa & \textbf{IA} & South Carolina & \textbf{SC} \\ \hline
Kansas & \textbf{KS} & South Dakota & \textbf{SD} \\ \hline
Kentucky & \textbf{KY} & Tennessee & \textbf{TN} \\ \hline
Louisiana & \textbf{LA} & Texas & \textbf{TX} \\ \hline
Maine & \textbf{ME} & Utah & \textbf{UT} \\ \hline
Maryland & \textbf{MD} & Vermont & \textbf{VT} \\ \hline
Massachusetts & \textbf{MA} & Virginia & \textbf{VA} \\ \hline
Michigan & \textbf{MI} & Washington & \textbf{WA} \\ \hline
Minnesota & \textbf{MN} & West Virginia & \textbf{WV} \\ \hline
Mississippi & \textbf{MS} & Wisconsin & \textbf{WI} \\ \hline
Missouri & \textbf{MO} & Wyoming & \textbf{WY} \\ \hline
\end{tabularx}
}\par
\sectiontitle{Canadian provinces and territories}
{\fontsize{9.8}{11.86}\selectfont
\rowcolors{2}{white}{RowTint}
\begin{tabularx}{\linewidth}{|Y{1.67}|Y{0.33}|Y{1.67}|Y{0.33}|}
\hline
\rowcolor{TableHead}
\textbf{Province / territory} & \textbf{Abbr.} & \textbf{Province / territory} & \textbf{Abbr.} \\ \hline
Alberta & \textbf{AB} & Nunavut & \textbf{NU} \\ \hline
British Columbia & \textbf{BC} & Ontario & \textbf{ON} \\ \hline
Manitoba & \textbf{MB} & Prince Edward Island & \textbf{PE} \\ \hline
New Brunswick & \textbf{NB} & Quebec & \textbf{QC} \\ \hline
Newfoundland and Labrador & \textbf{NL} & Saskatchewan & \textbf{SK} \\ \hline
Nova Scotia & \textbf{NS} & Yukon & \textbf{YT} \\ \hline
Northwest Territories & \textbf{NT} &  &  \\ \hline
\end{tabularx}
}\par
\end{document}
